\section{System Model and Actors}
\label{sec:model}

\paragraph{Devices.} A \emph{node} is a long-running process the user controls, on a computer that stays on, that holds keys, stores records, talks to other nodes, and watches the Bitcoin chain. A \emph{paired client} is the user-facing application talking to its owner's node: the native phone app (using embedded Tor or a sealed home-network channel), the desktop shell, or a browser on the node machine; in that mode it holds no spend keys and reaches other stores only through its node. Separately, a \emph{read-only public browser} (Tor Browser or an unpaired phone view) fetches signed catalogs directly from seller nodes and verifies every signature locally without running a node, but cannot place orders or pay (\S\ref{sec:public-browsing}). Finally, we specify a planned \emph{buyer-only light client} inside the signed desktop and mobile apps that derives deterministic buyer order and escrow keys from a 12-word recovery seed without running an always-on store node (\S\ref{sec:light-buyer}); this extension is designed but not yet shipped in the reference implementation.

\paragraph{Roles.}
\begin{description}[leftmargin=1.2em]
  \item[Buyer.] Browses the catalog, sends a sealed order, and pays. For moderated purchases, funds the escrow output and holds a presigned refund.
  \item[Vendor.] Publishes a signed store and listings, runs a node that opens sealed letters, ships, and collects payment. For moderated purchases, holds the presigned refund too, and collects from a stalling buyer through a dispute.
  \item[Moderator.] Publishes a record of its keys, filing fee, and acceptance policy, is listed in a store record, and is named in the signed escrow terms. Receives nothing unless a party opens a dispute. On a dispute, verifies a self-contained letter and reveals one secret scalar to the party it rules for, and publishes a signed record of the ruling and its outcome. Paid a fixed filing fee, outside the escrow, by the first party to file, before it decides.
  \item[Labeler.] Publishes signed sets of content labels about stores and listings. Buyers choose which labelers to follow; a labeler cannot change or remove anyone's records.
  \item[Replica.] A node that stores and serves catalog records it did not create, so that a store remains readable while its publisher is offline.
  \item[Directory operator.] Holds the directory key whose public half ships with the software, and signs the directory head: the introductions, the fee terms, and the key the fee is paid to (\S\ref{sec:directory}). Names no stores and cannot admit or remove one except by the fee rule every node checks.
  \item[Introduction.] A node with an onion address that the head lists, used to find stores. Introductions are untrusted for content: every record is verified against its signature, and every admission against the chain.
\end{description}

\paragraph{External parties.} The Bitcoin network and chain; a source of chain data for the node (its own Bitcoin node, compact block filters, or a block-explorer API); Tor relays; Lightning routing nodes and, for a vendor that uses one, a Lightning service provider; for stores that accept them, a Monero daemon and a Zcash compact-block server; an exchange-rate source for fiat estimates; and, for physical goods, a parcel carrier and, if a user opts in, a shipment-tracking service.

\paragraph{Records and messages.} A \emph{store record} names a store key, a display name, a region, direct-payment strings, the store's content rating, its mailbox key and onion address, the moderators it accepts, a profile, and the payment rails it accepts. An \emph{admission} ties a store key to a confirmed directory-fee payment. A \emph{listing} names its store key, describes one offer, and carries the vendor's content label. A \emph{label set} is a labeler's signed statement about stores and listings. An \emph{order letter} is an HPKE ciphertext addressed to one vendor. \emph{Escrow messages} are the three funding rounds and the release exchange between buyer and vendor. A \emph{dispute letter} goes from one party to the moderator, and a \emph{ruling} comes back.

\paragraph{Network assumptions.} We assume the Bitcoin chain makes progress and that a transaction paying a sufficient fee rate is confirmed within a bounded number of blocks; the escrow windows in \S\ref{sec:escrow} are sized against this assumption. We do not assume synchronous delivery between buyer and vendor nodes.
